\documentclass[11pt,a4paper]{article}
\usepackage[a4paper,margin=2.5cm]{geometry}
\usepackage[english]{babel}
\usepackage{amsmath}
\usepackage{amssymb}
\usepackage{amsthm}
\usepackage{array}
\usepackage{booktabs}
\usepackage{longtable}
\usepackage{hyperref}

\newtheorem{theorem}{Theorem}[section]
\newtheorem{lemma}[theorem]{Lemma}
\newtheorem{proposition}[theorem]{Proposition}
\newtheorem{conjecture}[theorem]{Conjecture}
\theoremstyle{definition}
\newtheorem{definition}[theorem]{Definition}
\newtheorem{remark}[theorem]{Remark}

\title{Computational Verification and 2D Tensor Network Analysis of Multi-Layer Crown Order Polytope Numerators}
\author{Wendel León Alvarado}
\date{\today}

\begin{document}

\maketitle

\begin{abstract}
We study Ehrhart $h^*$-polynomials of order polytopes associated with the multilayer crown posets $C_{k,n}$. A proposed reduction to a one-dimensional transfer trace is valid in height two but fails already for $C_{3,3}$: at dilation one, the exact lattice-point count is $28$, whereas the proposed trace gives $18$. We express the count as a contraction of a cylindrical two-dimensional tensor network with local bond dimension $D=m+1$. The contraction uses an optimized hierarchical order of multilinear tensor operations, summing shared indices according to Einstein convention, without boundary-rank truncation; the Ehrhart numerator is then recovered exactly from lattice-point counts by solving a triangular system in the binomial basis over $\mathbb{Z}$. Rational Sturm sequences certify simple, strictly negative real roots in 14 cases, including $C_{7,4}$ of degree 21. Exact symbolic reconstruction over $\mathbb{Q}[t]$ yields recurrences of order $n-1$ for the tested circumferences $n=3,4$. A universal recurrence of this order is formulated as a conjecture. These results constitute finite evidence, not a proof for all $k,n$.
\end{abstract}

\small\noindent\textbf{Keywords:} Crown posets, Ehrhart polynomials, tensor networks, Sturm method, real-rootedness.
\normalsize

\section{Introduction \& Background}

For a finite poset $P$, its order polytope is
\begin{equation}
\mathcal{O}(P)=\{x\in[0,1]^P : x_u\le x_v\text{ whenever }u\le_P v\}.
\end{equation}
It is a full-dimensional lattice polytope of dimension $|P|$; its vertices are the characteristic vectors of upper sets of $P$ \cite{Stanley1986}.

Fix integers $k\ge2$ and $n\ge3$. The poset $C_{k,n}$ has elements $x_{j,i}$ for $1\le j\le k$ and $i\in\mathbb{Z}/n\mathbb{Z}$. Its cover relations are
\begin{equation}
x_{j,i}\prec x_{j+1,i},\qquad x_{j,i}\prec x_{j+1,i+1}
\quad (1\le j<k),
\end{equation}
where the second index is read cyclically. The rank of $x_{j,i}$ is $j-1$, and $C_{k,n}$ has $d=kn$ elements.

Let
\begin{equation}
L_{k,n}(m)=|m\mathcal{O}(C_{k,n})\cap\mathbb{Z}^{d}|,
\qquad
\sum_{m\ge0}L_{k,n}(m)t^m=\frac{H_{k,n}(t)}{(1-t)^{d+1}}.
\end{equation}
Thus $H_{k,n}(t)=h^*(\mathcal{O}(C_{k,n}),t)$ and $H_{k,n}(0)=1$.

Lundström and Saud Maia Leite formulated the real-rootedness conjecture
for the simple one-dimensional crown family in work originating in 2025
\cite{LundstromLeite2026}; Cao and Fu subsequently proved the height-two
case $k=2$ \cite{CaoFu2026}. The present computational study addresses the
multilayer families $C_{k,n}$ with $k\ge3$. Its contributions are the
failure of the proposed one-dimensional transfer trace beyond height two,
a cylindrical two-dimensional tensor-network contraction, exact
reconstruction of the Ehrhart numerators, and finite-data recurrence
discovery. The recurrence and root-location statements for arbitrary
parameters remain conjectural.

\subsection{Methodological correction: the one-dimensional trace}

For $N\ge1$, define the matrix $K_N=(\min(a,b))_{1\le a,b\le N}$. In height two, setting $N=m+1$ gives the transfer identity
\begin{equation}\label{eq:height-two-trace}
L_{2,n}(m)=\operatorname{tr}(K_N^n).
\end{equation}
Indeed, fix the values on the bottom layer. For adjacent values $a_{i-1},a_i$, the top-layer value at their common upper neighbor has $N-\max(a_{i-1},a_i)$ choices. Replacing each bottom value $a_i$ by $N-a_i$ turns this factor into the matrix entry $\min(N-a_{i-1},N-a_i)$. Summing around the cycle yields \eqref{eq:height-two-trace}.

\subsection{Failure of the proposed higher-layer extension}

The height-two derivation sums out one layer and leaves pairwise interactions between neighboring values on a single cycle. With three or more layers, the intermediate variables remain coupled through two-dimensional local constraints; replacing those constraints by the proposed single transfer matrix does not preserve the count. The following exact example disproves that proposed extension.

\begin{proposition}\label{prop:transfer-counterexample}
For $C_{3,3}$ at dilation $m=1$, the lattice-point count is
\[
L_{3,3}(1)=28,
\]
whereas the proposed one-dimensional formula gives $\operatorname{tr}(K_2^3)=18$.
\end{proposition}

\begin{proof}
At $m=1$, represent each layer by the subset of its three cyclic positions assigned value $1$. A subset $A$ in one layer forces the next layer to contain $A\cup(A+1)$. Counting the compatible layer choices gives $18+6+3+1=28$. In contrast,
\[
K_2=\begin{pmatrix}1&1\\1&2\end{pmatrix},
\qquad \operatorname{tr}(K_2^3)=18.
\]
The two counts differ, so the proposed one-dimensional transfer-trace identity is not valid uniformly for $k\ge3$.
\end{proof}

\begin{remark}
This counterexample rules out the stated transfer-trace reduction; it does not rule out other one-dimensional representations for special parameters, nor does it disprove Conjecture~\ref{conj:realrooted}.
\end{remark}

\section{2D Boundary MPS Tensor Network Engine}

\subsection{Local factor and exact lattice-point count}

For a fixed dilation $m\ge0$, each coordinate takes values in $\{0,\ldots,m\}$. Define the local tensor
\begin{equation}\label{eq:local-tensor}
T_{a,b,c}=
\begin{cases}
1,&c\ge\max(a,b),\\
0,&\text{otherwise},
\end{cases}
\qquad 0\le a,b,c\le m.
\end{equation}
Its three indices have link dimension $D=m+1$. The indices $a,b$ encode adjacent values in one layer, and $c$ encodes the value above them. Thus the local factor enforces the required order constraint $x_{j+1,i}\ge\max(x_{j,i},x_{j,i+1})$.

The complete cylindrical network is the factor-graph contraction
\begin{equation}\label{eq:tensor-count}
L_{k,n}(m)=
\sum_{\substack{x_{j,i}\in\{0,\ldots,m\}\\
1\le j\le k,\ i\in\mathbb{Z}/n\mathbb{Z}}}
\prod_{j=1}^{k-1}\prod_{i\in\mathbb{Z}/n\mathbb{Z}}
T_{x_{j,i},\,x_{j,i+1},\,x_{j+1,i}}.
\end{equation}
The cyclic index creates the periodic boundary in the circumferential direction. Each coordinate is summed once, and every factor in the product enforces one local cover constraint.

\subsection{Boundary contraction and numerical modes}

The tensor network is contracted by organizing its multilinear factors
along the discrete cylindrical boundary and selecting an efficient
hierarchy of index contractions. This is an exact contraction of the
entire network: no boundary-rank compression or singular-value
truncation is applied. Accordingly, ``Boundary MPS'' refers here to the
boundary organization of the tensor contraction, not to an approximate
matrix-product-state method.

The exact contraction is carried out over arbitrary-precision integers,
so \eqref{eq:tensor-count} involves no floating-point rounding. For
spectral exploration, a separate numerical evaluation uses real or
complex double-precision tensor arithmetic on the selected computing
device. Numerical values are not used as integer certificates.

\section{Exact Algebraic Reconstruction \& Verification Results}

\subsection{Exact reconstruction of the Ehrhart numerator}

Put $d=kn$. Ehrhart theory gives
\begin{equation}\label{eq:binomial-basis}
L_{k,n}(m)=\sum_{j=0}^{d}h_j\binom{m+d-j}{d},
\qquad H_{k,n}(t)=\sum_{j=0}^{d}h_jt^j.
\end{equation}
We evaluate \eqref{eq:tensor-count} for $m=0,1,\ldots,d$ and solve the
resulting triangular linear system by exact integer arithmetic. The
coefficients above the structural degree $d-k$ vanish. No
floating-point interpolation is used.

\subsection{Exact results}

Table~\ref{tab:exact} reports the exact numerators obtained by tensor
contraction and binomial-basis reconstruction. All 14 tested cases
(100\%) have a rational Sturm certificate. In every case, the
Sturm count on $(-\infty,0)$ equals the degree and
$\gcd(H_{k,n},H'_{k,n})=1$, certifying that all roots are real, strictly
negative, and simple. The palindromic coefficients agree with the
reciprocity relation proved in Section~\ref{sec:structure}. The new boundary
case $C_{7,4}$ has degree $21$ and central coefficient
$56{,}806{,}162{,}004\approx5.68\times10^{10}$.

\small
\setlength{\tabcolsep}{3pt}
\begin{longtable}{@{}c c >{\raggedright\arraybackslash}p{0.49\textwidth} c c c@{}}
\caption{Complete exact dataset: $h^*$-polynomials, Sturm root counts, and JSON certificates.}\label{tab:exact}\\
\toprule
$k$ & $n$ & $H_{k,n}(t)$ & Degree & Negative roots & Certificate\\
\midrule
\endfirsthead
\caption[]{Complete exact dataset (continued).}\\
\toprule
$k$ & $n$ & $H_{k,n}(t)$ & Degree & Negative roots & Certificate\\
\midrule
\endhead
\midrule
\multicolumn{6}{r@{}}{\emph{Continued on next page}}\\
\endfoot
\bottomrule
\endlastfoot
2 & 3 &
$\begin{aligned}[t]
t^4 &{}+11t^3+24t^2+11t+1
\end{aligned}$ & 4 & 4/4 & Sturm + JSON\\
3 & 3 &
$\begin{aligned}[t]
t^6 &{}+18t^5+93t^4+160t^3\\
&{}+93t^2+18t+1
\end{aligned}$ & 6 & 6/6 & Sturm + JSON\\
4 & 3 &
$\begin{aligned}[t]
t^8 &{}+25t^7+211t^6+739t^5\\
&{}+1120t^4+739t^3+211t^2\\
&{}+25t+1
\end{aligned}$ & 8 & 8/8 & Sturm + JSON\\
3 & 4 &
$\begin{aligned}[t]
t^9 &{}+69t^8+1204t^7+7164t^6\\
&{}+16906t^5+16906t^4+7164t^3\\
&{}+1204t^2+69t+1
\end{aligned}$ & 9 & 9/9 & Sturm + JSON\\
5 & 3 &
$\begin{aligned}[t]
t^{10} &{}+32t^9+378t^8+2088t^7\\
&{}+5757t^6+8064t^5+5757t^4\\
&{}+2088t^3+378t^2+32t+1
\end{aligned}$ & 10 & 10/10 & Sturm + JSON\\
4 & 4 &
$\begin{aligned}[t]
t^{12} &{}+100t^{11}+3058t^{10}+37164t^9\\
&{}+204639t^8+553200t^7+767068t^6\\
&{}+553200t^5+204639t^4+37164t^3\\
&{}+3058t^2+100t+1
\end{aligned}$ & 12 & 12/12 & Sturm + JSON\\
6 & 3 &
$\begin{aligned}[t]
t^{12} &{}+39t^{11}+594t^{10}+4550t^9\\
&{}+19005t^8+44547t^7+59136t^6\\
&{}+44547t^5+19005t^4+4550t^3\\
&{}+594t^2+39t+1
\end{aligned}$ & 12 & 12/12 & Sturm + JSON\\
3 & 5 &
$\begin{aligned}[t]
t^{12} &{}+228t^{11}+11096t^{10}+175640t^9\\
&{}+1128655t^8+3314308t^7+4717744t^6\\
&{}+3314308t^5+1128655t^4+175640t^3\\
&{}+11096t^2+228t+1
\end{aligned}$ & 12 & 12/12 & Sturm + JSON\\
7 & 3 &
$\begin{aligned}[t]
t^{14} &{}+46t^{13}+859t^{12}+8468t^{11}\\
&{}+48263t^{10}+165086t^9+344061t^8\\
&{}+439296t^7+344061t^6+165086t^5\\
&{}+48263t^4+8468t^3+859t^2+46t+1
\end{aligned}$ & 14 & 14/14 & Sturm + JSON\\
5 & 4 &
$\begin{aligned}[t]
t^{15} &{}+131t^{14}+5869t^{13}+116355t^{12}\\
&{}+1144241t^{11}+6007079t^{10}+17690249t^9\\
&{}+30135467t^8+30135467t^7+17690249t^6\\
&{}+6007079t^5+1144241t^4+116355t^3\\
&{}+5869t^2+131t+1
\end{aligned}$ & 15 & 15/15 & Sturm + JSON\\
8 & 3 &
$\begin{aligned}[t]
t^{16} &{}+53t^{15}+1173t^{14}+14185t^{13}\\
&{}+103331t^{12}+473631t^{11}+1393927t^{10}\\
&{}+2657795t^9+3294720t^8+2657795t^7\\
&{}+1393927t^6+473631t^5+103331t^4\\
&{}+14185t^3+1173t^2+53t+1
\end{aligned}$ & 16 & 16/16 & Sturm + JSON\\
4 & 5 &
$\begin{aligned}[t]
t^{16} &{}+349t^{15}+32235t^{14}+1090390t^{13}\\
&{}+16336920t^{12}+122193552t^{11}\\
&{}+491902733t^{10}+1115855805t^9\\
&{}+1462319070t^8+1115855805t^7\\
&{}+491902733t^6+122193552t^5\\
&{}+16336920t^4+1090390t^3+32235t^2+349t+1
\end{aligned}$ & 16 & 16/16 & Sturm + JSON\\
6 & 4 &
$\begin{aligned}[t]
t^{18} &{}+162t^{17}+9641t^{16}+273808t^{15}\\
&{}+4129580t^{14}+35320040t^{13}\\
&{}+179517436t^{12}+560710448t^{11}\\
&{}+1100896046t^{10}+1376570220t^9\\
&{}+1100896046t^8+560710448t^7\\
&{}+179517436t^6+35320040t^5+4129580t^4\\
&{}+273808t^3+9641t^2+162t+1
\end{aligned}$ & 18 & 18/18 & Sturm + JSON\\
7 & 4 &
$ \begin{aligned}[t]
t^{21} &{}+193t^{20}+14374t^{19}+539238t^{18}\\
&{}+11319081t^{17}+141272409t^{16}\\
&{}+1094268784t^{15}+5432987952t^{14}\\
&{}+17714213382t^{13}+38591690550t^{12}\\
&{}+56806162004t^{11}+56806162004t^{10}\\
&{}+38591690550t^9+17714213382t^8\\
&{}+5432987952t^7+1094268784t^6\\
&{}+141272409t^5+11319081t^4+539238t^3\\
&{}+14374t^2+193t+1
\end{aligned}$ & 21 & 21/21 & Sturm + JSON\\
\end{longtable}
\normalsize

\subsection{Sanity Checks \& Data Verification}

The 14 certified polynomials were checked against the exact coefficient
lists in the JSON artifacts. Each satisfies
\[
H_{k,n}(0)=1,\qquad
t^{k(n-1)}H_{k,n}(1/t)=H_{k,n}(t),\qquad
\deg H_{k,n}=k(n-1).
\]
The second identity is strict palindromicity (the Gorenstein reciprocity
of these Ehrhart numerators); the degree agrees with
Lemma~\ref{lem:structure}. Independently, the JSON Sturm certificates
record a negative-root count equal to the degree and a square-free check
for each of the 14 cases. Thus these checks apply to the listed finite
dataset, not to the full infinite family.

\subsection{Rational Sturm certificates}

For a polynomial $H\in\mathbb{Z}[t]$, we form its Sturm sequence over $\mathbb{Q}$. Sign variations at $-\infty$ and at $0$ give the exact number of distinct roots in $(-\infty,0)$; this calculation involves rational polynomial arithmetic, not approximate root locations. Simplicity is checked independently by verifying $\gcd(H,H')=1$. If the negative-root count equals $\deg H$ and the gcd condition holds, all roots are real, simple, and strictly negative.

Sturm sequences can also isolate individual roots: for rational endpoints $a<b$, the difference in sign variations at $a$ and $b$ counts the roots in $(a,b)$, and rational bisection can refine an interval until it contains exactly one root. The current JSON output records the exact polynomial and Sturm-sequence coefficients, endpoint sign lists, variation counts, root count, and simplicity/reciprocity checks; it does not report decimal approximations as certificates.

Exact rational Sturm sequences and root isolations were computed for all
14 configurations, including $n=3$ for $k=2,\ldots,8$ and $n=4$ for
$k=3,\ldots,7$. Each execution produces an exact JSON certificate
recording the square-freeness verification, degree check, and endpoint
sign variations over $\mathbb{Q}$. The 14 exact certificates are
structured in JSON format, named
\texttt{certificado\_k\_n.json}, and available in the project's
supplementary repository.
\footnote{Execution scripts, CLI usage, and raw JSON certificate files are available in the open-source repository associated with this paper.}

\subsection{Degree and reciprocity}\label{sec:structure}

\begin{lemma}\label{lem:structure}
For $k\ge2$ and $n\ge3$, $H_{k,n}(t)$ is monic of degree $kn-k$ and satisfies
\begin{equation}\label{eq:reciprocity}
H_{k,n}(t)=t^{kn-k}H_{k,n}(1/t).
\end{equation}
\end{lemma}

\begin{proof}
Let $P=C_{k,n}$ and $d=kn$. Lattice points of $m\mathcal{O}(P)$ are weakly order-preserving maps $P\to\{0,1,\ldots,m\}$. Interior lattice points are maps $f:P\to\{1,\ldots,m-1\}$ that are strictly increasing on every cover relation.

Every maximal chain has $k$ elements. An interior lattice point therefore requires at least $k$ available integer values, so the smallest dilation with an interior lattice point is $m=k+1$. At this dilation, each vertical chain $x_{1,i}\prec\cdots\prec x_{k,i}$ forces $f(x_{j,i})=j$; these values satisfy all cover relations, giving exactly one interior lattice point. The codegree is consequently $k+1$. The standard degree-codegree relation for Ehrhart numerators gives $\deg H_{k,n}=d+1-(k+1)=d-k$, and the top coefficient equals the number of interior lattice points at the codegree, namely one. Thus $H_{k,n}$ is monic.

For reciprocity, let $\rho(x_{j,i})=j-1$. If $f$ is an interior lattice point of $m\mathcal{O}(P)$, then
\[
g(x)=f(x)-\rho(x)-1
\]
is weakly order-preserving and takes values in $\{0,\ldots,m-k-1\}$. Conversely, every such weak map gives an interior lattice point by $f(x)=g(x)+\rho(x)+1$. Hence the interior Ehrhart count is $E^\circ(m)=E(m-k-1)$ for $m\ge k+1$, and it is zero for smaller positive $m$. If $F(t)=\sum_{m\ge0}E(m)t^m$, this says
\[
\sum_{m\ge1}E^\circ(m)t^m=t^{k+1}F(t).
\]
Ehrhart reciprocity gives $\sum_{m\ge1}E^\circ(m)t^m=(-1)^{d+1}F(1/t)$. Substituting $F(t)=H_{k,n}(t)/(1-t)^{d+1}$ and simplifying yields \eqref{eq:reciprocity}.
\end{proof}

\begin{remark}\label{rem:minusone}
If $s=kn-k$ is odd, equation~\eqref{eq:reciprocity} forces $H_{k,n}(-1)=0$. Under Conjecture~\ref{conj:realrooted}, this root is simple and the other roots occur in reciprocal pairs.
\end{remark}

\begin{proposition}\label{prop:k23}
For $k=2$ and $n=3$,
\[
H_{2,3}(t)=t^4+11t^3+24t^2+11t+1,
\]
and its four roots are distinct and negative.
\end{proposition}

\begin{proof}
Equation~\eqref{eq:height-two-trace} gives $L_{2,3}(m)=1,18,129,571,1884$ for $m=0,1,2,3,4$. Substitution into \eqref{eq:binomial-basis} gives the displayed polynomial. Dividing by $t^2$ and setting $z=t+t^{-1}$ reduces its root equation to $z^2+11z+22=0$. Its roots $(-11\pm\sqrt{33})/2$ are distinct and less than $-2$, so each gives two distinct negative reciprocal roots of $t+t^{-1}=z$.
\end{proof}

\section{Symbolic Recurrence Reconstruction}
\label{sec:recurrences}

Solving exact coefficient systems over $\mathbb{Q}[t]$ against the
tensor-network dataset gives the following finite-range propositions.

\begin{proposition}[Finite exact recurrence for $n=3$]
\label{prop:recurrence-n3}
For the certified levels $k=2,\ldots,8$, the exact recurrence
\begin{equation}\label{eq:recurrence-n3}
H_{k,3}(t)
=(t^2+6t+1)H_{k-1,3}(t)
+(t^3-2t^2+t)H_{k-2,3}(t)
\end{equation}
holds for every tested transition $k=4,\ldots,8$. The recurrence
operator was reconstructed uniquely in the specified polynomial search
space, and the coefficient system is overdetermined by the exact data.
This proposition concerns only those certified transitions.
\end{proposition}

\begin{proof}
The coefficients are determined over $\mathbb{Q}$ from the certified
polynomials, then substituted into each transition using exact integer
polynomial arithmetic. The identity holds at $k=4,5,6,7,8$. The
uniqueness claim is relative to the bounded-degree search used by the
analyzer; it is not a uniqueness assertion among unrestricted operators
of arbitrary degree.
\end{proof}

\begin{proposition}[Finite exact recurrence for $n=4$]
\label{prop:recurrence-n4}
For the certified levels $k=3,\ldots,7$, the exact recurrence
\begin{equation}\label{eq:recurrence-n4}
H_{k,4}(t)
=P_4(t)H_{k-1,4}(t)
+Q_4(t)H_{k-2,4}(t)
+R_4(t)H_{k-3,4}(t)
\end{equation}
holds for the tested transitions $k=6,7$, where
\begin{align*}
P_4(t)&=t^3+23t^2+23t+1,\\
Q_4(t)&=7t^5-28t^4-22t^3-28t^2+7t,\\
R_4(t)&=t^8-5t^7+9t^6-5t^5\\
&\qquad{}-5t^4+9t^3-5t^2+t.
\end{align*}
Within the specified degree bounds and reciprocal-polynomial ansatz,
the exact coefficient system is overdetermined and has a unique integral
solution. This proposition asserts the identity only at $k=6,7$.
\end{proposition}

\begin{proof}
The three operators were reconstructed over $\mathbb{Q}$ using reciprocal
spans $(3,6,9)$ and degree bounds $(4,7,10)$. Exact substitution into
$H_{6,4}$ and $H_{7,4}$ gives zero residual in both transitions. The
uniqueness statement is relative to this finite-dimensional ansatz.
\end{proof}

The recurrence operators were uniquely reconstructed over
$\mathbb{Q}[t]$ and certified on exact integer arithmetic. Here
``uniquely'' is relative to the bounded-degree search space for $n=3$
and, for $n=4$, to the stated degree bounds and reciprocal spans; it does
not assert uniqueness without these restrictions.

All operator coefficients lie in $\mathbb{Z}[t]$ and are reciprocal with
their indicated spans:
\[
\begin{aligned}
P_3(t)&=t^2P_3(1/t),&
Q_3(t)&=t^4Q_3(1/t),\\
P_4(t)&=t^3P_4(1/t),&
Q_4(t)&=t^6Q_4(1/t),&
R_4(t)&=t^9R_4(1/t),
\end{aligned}
\]
where $P_3(t)=t^2+6t+1$ and $Q_3(t)=t^3-2t^2+t$. The order is $n-1$
for the two circumferences investigated. The $n=3$ identities were
checked for transitions $k=4,\ldots,8$. For $n=4$, the two independent
transitions $k=6,7$ determine a unique integral solution under reciprocal
spans $(3,6,9)$ and degree bounds $(4,7,10)$.

\begin{remark}
The coefficients shown here are obtained by exact reconstruction and
substitution into the certified polynomials. In particular, the
$t^3$ coefficient of $Q_3$ and the $t^8$ coefficient of $R_4$ are both
$+1$; changing either to $-1$ does not reproduce the dataset. These
computations establish the identities only on the stated finite ranges.
\end{remark}

\section{Real-Rootedness \& Interlacing Implications}
\label{sec:root-implications}

\begin{conjecture}[Higher-layer negative real-rootedness]
\label{conj:realrooted}
For every $k\ge3$ and $n\ge3$, all zeros of $H_{k,n}(t)$ are negative
real numbers and are simple.
\end{conjecture}

\begin{conjecture}[Cylindrical recurrence of order $n-1$]
\label{conj:universal-recurrence}
For each fixed circumference $n\ge3$, there are polynomials
$A_{1,n}(t),\ldots,A_{n-1,n}(t)\in\mathbb{Z}[t]$, independent of $k$,
such that the recurrence holds for every $k\ge n$, where all lagged
polynomials $H_{k-j,n}$ are defined:
\[
H_{k,n}(t)=\sum_{j=1}^{n-1}A_{j,n}(t)H_{k-j,n}(t).
\]
\end{conjecture}

The recurrence discovery in Section~\ref{sec:recurrences} is consistent
with this order-$n-1$ conjecture for $n=3,4$, but its exact identities are
verified only on the finite transition ranges stated there. Neither
finite recurrence data nor real-rootedness of the listed examples proves
Conjecture~\ref{conj:universal-recurrence}.

\subsection{Characteristic discriminant for the $n=3$ recurrence}

With the verified convention $H_{k,3}=P_3H_{k-1,3}+Q_3H_{k-2,3}$, the
characteristic equation in the layer index is
$\lambda^2-P_3(t)\lambda-Q_3(t)=0$. Its discriminant is therefore
\[
\Delta_3(t)=P_3(t)^2+4Q_3(t)
=t^4+16t^3+30t^2+16t+1.
\]
For $t<0$, write $x=-t>0$ and $z=x+x^{-1}\ge2$. Then
\[
\frac{\Delta_3(-x)}{x^2}=z^2-16z+28=(z-2)(z-14).
\]
Consequently, $\Delta_3(t)$ vanishes at $z=2$ and $z=14$, is positive
when $z>14$, and is negative when $2<z<14$. Thus the characteristic
multipliers are not real for every negative $t$; this discriminant does
not by itself imply stability or real-rootedness of the polynomial
sequence. The Sturm certificates establish the latter only for the
individual cases in Table~\ref{tab:exact}.

The alternative expression
$P_3(t)^2-4(-t^3-2t^2+t)=t^4+16t^3+46t^2+8t+1$
is not positive for every $t<0$: at $t=-4$ it equals $-63$. Moreover,
the proposed $Q_3(t)=-t^3-2t^2+t$ differs from the verified operator
$t^3-2t^2+t$ by $-2t^3$, so it does not reproduce the certified
recurrences. Similarly, changing the leading coefficient of $R_4$ from
$+1$ to $-1$ changes the recurrence residual by a nonzero multiple of
$2t^8H_{k-3,4}$. The coefficients in Section~\ref{sec:recurrences} are
retained because they reproduce the exact data.

\subsection{Interlacing status}

The Sturm certificates prove that each listed polynomial has only simple
negative real roots. They do not establish strict interlacing between
different values of $k$. The degrees increase by $n-1$ in each fixed-$n$
family, so any interlacing claim requires a separately specified
generalized interlacing criterion and an independent proof; no such
general theorem is inferred from the current finite dataset.

\section{Concluding Remarks \& Future Work}

The height-two transfer trace is justified by summing out a single layer, but the explicit $C_{3,3}$ count shows that the proposed one-dimensional extension does not preserve lattice-point counts for higher layers. The cylindrical tensor network retains the full local coupling, while exact contraction and binomial-basis reconstruction yield reproducible integer polynomials. Sturm certificates verify simple negative real-rootedness for the 14 finite cases in Table~\ref{tab:exact}. The recurrence identities are verified only on the displayed finite ranges, and the universal order-$n-1$ recurrence and higher-layer real-rootedness remain conjectural. Future work includes deriving recurrences from the tensor-network transfer structure and proving root-location or generalized interlacing results uniformly in $k,n$.

\subsection{Data \& Code Availability}

The project directory contains the tensor-network engine and exact
reconstruction modules in the \texttt{core/} directory, the command-line
entry point \texttt{main.py}, and the 14 JSON certificates in
\texttt{certificados\_json/}. These paths are relative to the project
subdirectory \texttt{ARQUITECTURA CON TENSORES/}. Sturm verification is
carried out by \texttt{core/root\_verifier.py}; recurrence and
interlacing computations are carried out by
\texttt{core/interlacing\_analyzer.py}. The certificates record the
polynomial coefficients and exact Sturm data needed to reproduce the
finite checks reported here.

\begin{thebibliography}{9}
\bibitem{LundstromLeite2026}
T.~Lundstr\"om and L.~Saud Maia Leite,
\emph{Order polytopes of crown posets},
European Journal of Combinatorics \textbf{133} (2026), 104304.
\href{https://doi.org/10.1016/j.ejc.2025.104304}{doi:10.1016/j.ejc.2025.104304}.

\bibitem{Stanley1986}
R.~P.~Stanley,
\emph{Two poset polytopes},
Discrete \& Computational Geometry \textbf{1} (1986), 9--23.

\bibitem{CaoFu2026}
Y.~Cao and H.~Fu,
\emph{Real-rooted Ehrhart $h^*$-polynomials of crown posets and wheel and whirl matroids},
preprint (2026), arXiv:2610.40496.
\href{https://arxiv.org/abs/2610.40496}{arXiv:2610.40496}.
\end{thebibliography}

\end{document}
